problem:  
Let \( (S^3,g_\lambda) \) be the Berger family of metrics on the 3‑sphere, defined by scaling the fibers of the Hopf fibration  
\[
S^1 \longrightarrow S^3 \xrightarrow{\pi} S^2
\]
by a factor \(\lambda>0\).  As \(\lambda \to 0\), the metrics \(g_\lambda\) collapse \(S^3\) to \(S^2\) with uniformly bounded sectional curvature.  
Let \(A_\lambda\) denote the connection 1‑form on the Hopf bundle defined by the orthogonal complement of the fibers, and let \(F_\lambda=dA_\lambda\) denote its curvature 2‑form on \(S^2\).  

**Conjecture (Curvature Rigidity under Ricci‑bounded Collapse).**  
Suppose \( (S^3,g'_\lambda) \) is any smooth one‑parameter family of Riemannian metrics such that  
1. each \( (S^3,g'_\lambda) \) realizes a Riemannian submersion \( \pi_\lambda:S^3\to S^2\) with \(S^1\) fibers;  
2. the fiber length tends uniformly to zero as \(\lambda\to 0\);  
3. the Ricci curvature of \(g'_\lambda\) remains uniformly bounded.  

Then the curvature densities of the induced connection forms \(F'_\lambda=dA'_\lambda\) on \(S^2\) satisfy  
\[
F'_\lambda \xrightarrow{\;\text{weakly}\;} c\,\omega_{S^2}
\]
for a constant \(c\neq 0\) determined by the Euler class of the bundle.  
Equivalently, **bounded‑Ricci collapse of \(S^3\) with circle fibers cannot produce measure‑valued curvature concentration:** the limiting curvature distribution is always absolutely continuous with respect to the smooth area measure on \(S^2\).  

Why is it a "good" problem:  
This formulation turns a broad geometric idea into a precise conjecture about rigidity of curvature distributions under collapse.  It is mathematically consistent, touches central open issues in **Ricci‑bounded collapse, gauge theory, and geometric analysis**, and can be tested using tools from Cheeger–Colding limit theory, geometric measure theory, and Yang–Mills analysis.  A proof or counterexample would clarify whether small‑fiber collapse preserves smooth curvature distribution or allows microscopic “bubbling.”  The statement is simple—just the Hopf fibration, Ricci bounds, and weak convergence—but resolving it would yield deep new insight into the analytic structure of collapsing manifolds, illustrating the elegance, depth, and cross‑disciplinary nature of a truly "good" mathematical problem.